\section{Related work}
\label{sec:related}

Four literatures meet in this benchmark: coloured-noise modelling, model-free
decomposition, machine learning on coordinate series, and shared evaluation
protocols. Table~\ref{tab:related} summarises what each provides and what it
leaves out; the text keeps only the points that shape the protocol.

\paragraph{Coloured noise and the trajectory model.}
Power-law noise in daily coordinates and its control on velocity precision are
the founding results of the field \citep{mao1999noise,williams2003coloured}.
The maximum-likelihood trajectory fit, with efficient covariance construction and
support for missing epochs, is implemented in CATS \citep{williams2008cats} and
Hector \citep{bos2008fast,bos2013fast}, and the noise model has been extended to
multivariate and seasonal formulations
\citep{amirisimkooei2007noise,amirisimkooei2009multivariate,klos2018seasonal},
including stochastic seasonality in velocity estimation
\citep{xu2025velocity}.
This literature supplies our reference estimator and the generalised-least-squares
velocity covariance. It treats the trajectory model as the frame within which
everything is measured; we instead score it as one candidate among several,
including on gap filling and detection, which its formulation was not designed
for.

\paragraph{Model-free decomposition.}
Singular spectrum analysis \citep{wdowinski1997spatial,ji2024ssa}, wavelet
shrinkage \citep{ji2020wavelet}, multivariate SSA
\citep{gruszczynska2017mssa} and spatiotemporal PCA/Karhunen--Lo\`eve filtering
\citep{dong2006spatiotemporal} separate seasonal signal from residual without a
parametric trajectory. They are
the natural classical competitors for a learned model, and we carry SSA, wavelet
shrinkage and a moving median as three independent members of this family. Their
usual validation is the task they were built for, on the series they were built
on; here they are scored on the same injected realisations as every other method.

\paragraph{Learned processing.}
Networks have been applied to discontinuity detection
\citep{crocetti2021discontinuity,wang2025changepointcnn}, denoising
\citep{li2018denoising,mastella2025denoising,gao2022gbdt,he2026tspredictor},
gap filling as low-rank completion \citep{liu2022imputation}, network-wide
denoising with graph models \citep{bachelot2025gnn}, and uncertainty propagation
\citep{kiani2022uncertainty}, and decomposition-plus-learning pipelines have been
evaluated for height-series prediction \citep{guo2026decomposition}. Time-domain
machine learning for GNSS series has also been reviewed as a whole
\citep{gabsatarov2025review}; related variants continue to appear
\citep{simsek2025dl,ozarpaci2024machine}. The evidence is usually recovery error
on the task the network was trained for. Where a geodetic product is reached, the
velocity is often estimated directly from positions rather than through the
coloured-noise model, so the effect of processing on the published uncertainty is
not traced. That is the step this benchmark adds.

\paragraph{Shared protocols.}
DOGEx fixed a detection dataset, injected known offsets and had independent
groups run detectors blind \citep{gazeaux2013detecting}; it remains the model for
a community protocol, but it covers one task and synthetic truth. CATS and Hector
are tools, not protocols. We are not aware of a protocol that fixes the data and
the split, creates exact truth by injection, scores classical and learned methods
on the same realisations across three tasks, validates detection against real
catalogue labels, and carries each output through to velocity, noise amplitudes
and interval coverage.

\begin{table*}[t]
\centering
\caption{Positioning of representative work. ``Downstream'' means the study
reports the effect of processing on a geodetic product (velocity or its
uncertainty); ``common truth'' means classical and learned methods are scored on
the same known realisation.}
\label{tab:related}
\footnotesize
\setlength{\tabcolsep}{3pt}
\begin{tabularx}{\textwidth}{@{}p{2.4cm}p{2.0cm}Xp{1.7cm}p{2.8cm}@{}}
\toprule
Study or family & Task & Evidence base & Downstream & What this benchmark adds \\
\midrule
MLE trajectory; CATS, Hector & noise, velocity & parametric fit to real series & yes & scores the same estimator against nonparametric truth and other methods \\
SSA, wavelet, mSSA & decomposition & recovery on selected series & no & common realisations, three tasks, seed-aware learned comparison \\
Detection networks & offset detection & synthetic or station-specific labels & no & real catalogue labels, amplitude and SNR strata, delay \\
Denoising, imputation, GNN networks & T1, T2 & recovery error & no & downstream noise and uncertainty, gap lengths to one year \\
Learned uncertainty and velocity & product & direct velocity from positions & partial & processing-to-noise-to-uncertainty chain \\
DOGEx & detection & blind synthetic injection & no & three tasks, real labels, two analysis centres \\
This paper & T1--T3 & injection, two truth families, real labels & yes & station-level split, 10--20 seeds, two architectures, RW arm \\
\bottomrule
\end{tabularx}
\end{table*}
